\begin{center}
	\zihao{4}\heiti
	%
	\LARGE{ 中国科学院大学\\$\boldsymbol{2016}$年招收攻读硕士学位研究生入学统一考试试卷\\
	科目名称：电动力学}

\end{center}
\noindent
考生须知：\\
1.本试卷满分为150分,全部考试时间总计180分钟。\\
2.所有答案必须写在答题纸上,写在试题纸上或草稿纸上一律无效。

\vspace*{2ex}
\hrule
\vspace*{4ex}

\fancypagestyle{mypagestyle}{%
	\fancyhf{}% Clear header/footer
	\fancyfoot[OR]{第~\thepage~页，共2页}%
	\fancyfoot[EL]{第~\thepage~页，共2页}%
	\fancyfoot[OL]{科目名称：电动力学 }%
	\fancyfoot[ER]{科目名称：电动力学 }%
	\renewcommand{\headrulewidth}{0pt}% Header rule of .4pt
	\renewcommand{\footrulewidth}{0.5pt}% Header rule of .4pt
}
\fancypagestyle{plain}{%
	\fancyhf{}% Clear header/footer
	\fancyfoot[OR]{第~\thepage~页，共2页}%
	\fancyfoot[EL]{第~\thepage~页，共2页}%
	\fancyfoot[OL]{科目名称：电动力学 }%
	\fancyfoot[ER]{科目名称：电动力学 }%
}

\setcounter{page}{1}%设置页面计数器
\pagestyle{mypagestyle}

\onehalfspacing
% 排版结束，下面是题目

\vspace{0.4cm}

\noindent
一、简答题（共36分，每题12分）
\vspace{0.2cm}

\noindent
1.\hspace{0.8em}写出真空无源情况下的麦克斯韦方程组，由此写出电场强度，磁场强度各自满足的波动方程。\vspace{0.3cm}

\noindent
2.\hspace{0.8em}写出导体复电容率的表达式，解释其实部、虚部的物理含义，并由此写出良导体的条件。\vspace{0.3cm}

\noindent
3.\hspace{0.8em}写出洛伦兹规范下真空电磁场的标势、矢势满足的波动方程。对于一般的时变电荷分布，电流分布，写出该波动方程的推迟势解。\vspace{1cm}

\noindent
二、\hspace{0.2em}单项选择题（共36分，每题6分）
\vspace{0.2cm}

\noindent
1.\hspace{0.8em}已知电介质区域内的自由电荷分布。将空间中静电场，静电势分别记为$\Vec{E}$，$ \phi$。给定下列哪个边界条件无法使区域内的电场唯一确定:\vspace{0.2cm}

A.\hspace{0.8em}边界上$ \phi$\hspace{11.2em}
B.\hspace{0.8em}边界上$\Vec{E_n}$（法向分量）\vspace{0.2cm}

C.\hspace{0.8em}边界上$\Vec{E_t}$（切向分量）\hspace{5em}
D.\hspace{0.8em}边界上部分区域$ \phi$，其余区域$\Vec{E_n}$\vspace{0.2cm}

\noindent
2.\hspace{0.8em}正四面体的各棱均匀分布着电荷。取四面体中心为坐标原点。下列说法正确的是：\vspace{0.2cm}

A.\hspace{0.8em}电偶极距，电四极距均为零\hspace{6em}
B.\hspace{0.8em}电偶极距为零，电四极距不为零\vspace{0.2cm}

C.\hspace{0.8em}电偶极距，电四极距均不为零\hspace{5em}
D.\hspace{0.8em}电偶极距不为零，电四极距为零\vspace{0.2cm}

\noindent
3.\hspace{0.8em}真空矩形波导管边长分别是a、b( a > b )，该波导管能传播的最大波长为：\vspace{0.2cm}

A.\hspace{0.8em}2a\hspace{5em} B.\hspace{0.8em}2b\hspace{5em}  C.\hspace{0.8em}b\hspace{5em} D.\hspace{0.8em}4a\vspace{0.2cm}

\noindent
4.\hspace{0.8em}频率为$f$的电磁波入射到介电常数为$\varepsilon$、磁导率为$\mu$、电导率为$\sigma$的良导体表面，其穿透深度的表达式为\vspace{0.3cm}

A.\hspace{0.8em}$\displaystyle{ \sqrt{\frac{2}{f \mu \varepsilon} } }$\hspace{5em}
B.\hspace{0.8em}$\displaystyle{ \sqrt{\frac{1}{\pi  f \mu \varepsilon} } }$\hspace{5em}
C.\hspace{0.8em}$\displaystyle{ \sqrt{\frac{2}{f \varepsilon \sigma} } }$\hspace{5em}
D.\hspace{0.8em}$\displaystyle{ \sqrt{\frac{ 2f }{ \mu \sigma} } }$\vspace{0.3cm}

\noindent
5.\hspace{0.8em}对于电偶极辐射电磁场，下列说法正确的是（取偶极方向为极轴）：\vspace{0.2cm}

A.\hspace{0.8em}电场强度反比于距离，磁感应强度反比于距离的平方\vspace{0.1cm}

B.\hspace{0.8em}辐射总功率正比于距离\vspace{0.1cm}

C.\hspace{0.8em}辐射能流密度反比于距离的平方\vspace{0.1cm}

D.\hspace{0.8em}电场强度沿方向角方向，磁感应强度沿极角方向

\vspace{0.3cm}

\noindent
6.\hspace{0.8em}下面哪个量在规范变换和洛伦兹变换下都是不变的\vspace{0.3cm}

A.\hspace{0.8em} $\Vec{A}^2$\hspace{5em} B.\hspace{0.8em} $\Vec{A} \cdot \Vec{J} - \varphi \rho$ \hspace{5em} C.\hspace{0.8em} $\Vec{J}$ \hspace{5em} D.\hspace{0.8em} $\Vec{E}^2 - c^2 \Vec{B}^2 $

\newpage

\noindent
三、\hspace{0.2em}（共30分）真空中放置一个半径为$a$的导体球，球心位于原点。在过球心的z轴（极轴）上，在球外$z = -R$处放置电荷$+Q$，在$z = +R$处放置电荷$-Q$
\vspace{0.2cm}

\noindent
1.写出电势在球面处满足的边界条件；

\noindent
2.用电像法求解导体球外距原点为$r$，极角为$\theta$处的电势;

\noindent
3.考虑$a \ll R$，$r \ll R$的情况，将电势表达式展开保留到$1/R^2$次项


\noindent
4.求上述情况下导体球表面的感应电荷密度以及总电荷


\noindent
5.求导体球的电偶极距
\vspace{1cm}

\noindent
四、\hspace{0.2em}（共30分）有一个半径为$r_0$，磁导率为$\nu$的电介质球被永久磁化。取球心为球坐标原点，\vspace{0.2cm}

\noindent
其极化矢量可表达为$\Vec{P}= \displaystyle{ \frac{P_0}{r^2} } \Vec{e_r} $，其中$P_0$为常数。今使电介质球绕极轴以角速度$\omega$转动。
\vspace{0.2cm}

\noindent
1.写出球内非球心处的束缚电荷密度和球面上的电荷密度;

\noindent
2.用磁标势法求球内外的磁场；

\noindent
3.求出球内的磁化强度，磁化电流密度以及球面上的磁化电流密度。\vspace{0.2cm}

\noindent
[\quad 提示：球坐标系中$\nabla \cdot \Vec{f} = \displaystyle
{ 
\frac{1}{r^2}} \frac{\partial}{\partial r} \left(  r^2 f_r\right)
+ 
\frac{1}{r \sin \theta} \frac{\partial}{\partial \theta} \left(  \sin \theta f_{\theta} \right) 
+
\frac{1}{r \sin \theta}\frac{\partial f_{\phi} }{ \partial \phi } ,
\\[2mm]
\vspace{1cm}
\nabla \varphi = \frac{\partial \varphi }{ \partial r} \Vec{e_r} + \frac{1}{r} \frac{ \partial \varphi }{ \partial \theta} \Vec{e_{\theta} } + \frac{1}{r \sin \theta} \frac{\partial \varphi}{ \partial \phi} \Vec{e_{\phi} }
$ \quad]

\vspace{0.5cm}
\noindent
五、\hspace{0.2em}（共18分）
如图所示，厚度为a的介质薄层夹在介质1和介质2之间（介质1，2可视为无限厚），它们的介电常数分别为$\varepsilon$，$\varepsilon_1$，$\varepsilon_2$且满足$\varepsilon > \varepsilon_1 > \varepsilon_2$，磁导率均为$\mu_0$。现有一平面电磁波入射到介质薄层，入射方向与界面夹角为$\varphi$，电场偏振方向垂直于它在界面上的入射面。为保证电磁波能在薄层中以稳定的模式传播，其能量不从介质1和介质2中透射出去，需要满足如下两个
条件：（1）在界面上全反射；（2）满足驻波条件$2a \sqrt{\varepsilon / \varepsilon_0 } \displaystyle{\frac{\omega}{c}}  \sin \varphi + \phi_1 + \phi_2 = 2n \pi , \quad n = 0, 1, 2, 3...$\vspace{0.1cm}

\noindent
其中$\phi_1, \phi_2$分别是电磁波在两个界面上全反射时，反射波与入射波之间的相位差。\vspace{0.2cm}

\noindent
1.$\varphi$满足什么条件才能保证电磁波在两个界面上全反射？\vspace{0.2cm}

\noindent
2.已知菲涅尔公式$\displaystyle{ \frac{ E^{'} } {E} = - \frac{ \sin (\theta - \theta^{''} ) }{ \sin (\theta + \theta^{''} )  }   }$，$E, E^{'}$分别为入射电磁波和反射电磁波的复振幅，

\noindent
$\theta, \theta^{''} $分别为入射角和折射角。求电磁波分别在两个界面上全反射时的相位差$\phi_1, \phi_2$;
[提示: $\displaystyle{\frac{E^{'}}{E} = e^{i \theta} } $]\vspace{0.2cm}

\noindent
3.若$\varepsilon_1 = \varepsilon_2, \cos = \sqrt{\varepsilon_1 / \varepsilon}$，求传播波的截止频率（非零解）

\vspace{-0.8cm}
\begin{figure}[h!]
    \flushright
    \includegraphics[height=7cm]{test/2016/pic_1.png}
\end{figure}
